\begin{abstract}
\noindent
Richard Nixon's secret White House taping system recorded 2,719 hours of conversation between February 16, 1971 and July 18, 1973. Hand transcription reached only the conversations needed for trials, diplomatic histories and edited selections. The Nixon Corpus transcribes all 19,426 conversations released as individual recordings: 2,206 hours, 19.0 million words, every word timed and every speech segment labeled. It rests on the first publicly documented speech recognizer fine-tuned on the tapes, a Transformer model (NVIDIA's Parakeet TDT 0.6B), and a second, independent recognizer (Whisper large-v3), reconciled by a large language model against the National Archives' subject logs, H.\,R. Haldeman's diary and the day's other conversations. Every change is a reversible edit. A voiceprint classifier names fourteen principal speakers. On 165 held-out conversations transcribed by the State Department's historians, the corpus recovers 76.1\% of words and 80.1\% of names; Whisper alone recovers 69.2\% and 60.0\%. On the 130 of them outside the classifier's training data, 92.4\% of named words carry the right speaker. Transcripts, model and code are open.
\end{abstract}

\noindent\textbf{Keywords:} Nixon White House tapes; presidential recordings; archives; speech recognition; speaker identification; large language models; Watergate
